\customchapter{ABSTRACT}

\begin{center}
\large \vspace{-1.5cm} \textbf{HyperNCD--Sonata hybrid model for novel class
discovery in the semantic segmentation of architectural heritage point clouds}
\end{center}

Laser scanners used to record heritage buildings produce point clouds with
millions of points, but those points come with no meaning attached. The file
knows where matter is, not what element it belongs to. Someone still has to go
through the model by hand and label the columns, the arches, the vaults, the
cornices and the buttresses one at a time, and there is no tool that does this
automatically. Existing 3D segmentation models do not help either. They are
trained on a fixed list of classes, so when a point does not belong to any of
them the model does not report that it is unsure. It simply assigns the closest
class it knows. In a heritage building this is a serious problem, because the
elements worth finding are precisely the ones that no catalogue lists. Our goal
is to improve, and to actually measure, how well those unlisted elements are
recovered.

Novel class discovery does allow a model to group classes it has never seen
labelled, and its best current method is HyperNCD. Its results on those classes,
however, are weak: between 11.8 and 47.5 mIoU, while an equivalent supervised
model reaches 49.8. When we looked into why, we found that the method builds its
prototypes from geometric descriptors computed by hand, such as linearity,
planarity and scattering. A more recent work, Sonata, shows with numbers that
relying on this kind of low-level signal damages the representation, and that
fixing it raises the score of a frozen descriptor from 21.8 to 72.5 mIoU. We
therefore propose a hybrid model that replaces those hand-made descriptors with
the self-supervised encoder of Sonata inside the prototype module of HyperNCD.
Everything else stays as it is, so if the metric moves we know the replacement
caused it. We evaluate the model on a real survey of the Cathedral of Lima, about
two gigabytes after voxelisation, split into five separate sections.

We first reproduce the original method over three independent runs, so that we
know where we are starting from. Against that reference we measure how much the
hybrid gains on the novel classes and check that the known classes do not get
worse. We then compare three versions of the descriptor: the hand-made one, the
self-supervised one, and both of them together. Standard metrics are not enough
here, because a hybrid inherits the strengths and the failures of the two methods
it combines, so we also define a closed-gap measure that tells us what share of
the distance between known and novel classes the model manages to recover.
Finally, we report the metric halfway through the pipeline as well as at the end,
which lets us say whether an improvement came from the representation or from the
hypergraph reasoning.

The question behind all this is whether a weakness proven in a setting with full
supervision is still a weakness when there are no labels for the target classes,
where the representation is all the model has to rely on. As far as we know, no
previous work has measured that. In practical terms, every extra point of mIoU is
manual correction that the engineer no longer has to do. Specialists are already
satisfied with the accuracy they get by hand, so the point is not to beat them on
accuracy but to automate the task well enough to be worth using. \\

\noindent \textbf{Keywords:}\\
\noindent Point clouds; 3D semantic segmentation; Novel class discovery;
Self-supervised learning; Hypergraphs; Architectural heritage
